(1) $|\uparrow\uparrow\rangle$ 与 $|\downarrow\downarrow\rangle$ 的交换项均为 $A/4$，Zeeman项分别为 $\pm(\gamma_e+\gamma_n)B/2$，所以 $E_{\pm1}=A/4\pm(\gamma_e+\gamma_n)B/2$；这些态不与 $m_F=0$ 混合。
(2) 在 $(|\uparrow\downarrow\rangle,|\downarrow\uparrow\rangle)$ 基中，$H_0=\begin{pmatrix}-A/4+\delta\gamma B/2&A/2\\A/2&-A/4-\delta\gamma B/2\end{pmatrix}$，$\delta\gamma=\gamma_e-\gamma_n$。本征能量 $E_{0,\pm}=-A/4\pm\frac12\sqrt{A^2+(\delta\gamma B)^2}$。
(3) $B=0$ 时三重态 $F=1$ 的三个态均为 $A/4$，单重态 $F=0$ 为 $-3A/4$。若 $|\delta\gamma B|\ll A$，$E_{0,+}=A/4+(\delta\gamma B)^2/(4A)+O(B^4)$，$E_{0,-}=-3A/4-(\delta\gamma B)^2/(4A)+O(B^4)$。
(4) 当 $|\delta\gamma B|\gg A$，$m_F=0$ 两态趋向Zeeman积态，能量近似 $-A/4\pm|\delta\gamma B|/2$；有限交换耦合使两态最小能隙为 $A$，在 $B=0$ 避免交叉。$\gamma B$ 和 $A$ 都是能量，平方根内各项量纲一致。
